% ============================================================================
%  Bibliografia (kolejność według pierwszego cytowania)
% ============================================================================
\newcommand{\dostep}{[Dostęp: 11.11.2026]}

\begingroup
\small
\setstretch{1.0}
\begin{thebibliography}{99}
\phantomsection
\addcontentsline{toc}{chapter}{Bibliografia}
\setlength{\itemsep}{2pt}

\bibitem{unity-order} Unity Technologies, „Order of execution for event functions,” Unity
  Manual 2022.3. [Online]. Dostępne:
  \url{https://docs.unity3d.com/2022.3/Documentation/Manual/ExecutionOrder.html}. \dostep.

\bibitem{r3} Cysharp, „R3 -- The new future of dotnet/reactive and UniRx,” repozytorium
  GitHub. [Online]. Dostępne: \url{https://github.com/Cysharp/R3}. \dostep.

\bibitem{unirx} Y. Kawai, „UniRx -- Reactive Extensions for Unity,” repozytorium GitHub
  (zarchiwizowane 16.02.2024). [Online]. Dostępne: \url{https://github.com/neuecc/UniRx}.
  \dostep.

\bibitem{gra-smb} „Super Monkey Ball (video game),” Wikipedia. [Online]. Dostępne:
  \url{https://en.wikipedia.org/wiki/Super_Monkey_Ball_(video_game)}. \dostep.

\bibitem{gra-miu} The Marble Collective, „Marble It Up!,” oficjalna strona gry. [Online].
  Dostępne: \url{https://marbleitup.com/}. \dostep.

\bibitem{gra-golf} Miniclip, „Golf Battle,” Google Play. [Online]. Dostępne:
  \url{https://play.google.com/store/apps/details?id=games.onebutton.golfbattle}. \dostep.

\bibitem{gra-golf-recenzja} TapSmart, „Golf Battle: fast-paced online mini golf,”
  21.12.2019. [Online]. Dostępne:
  \url{https://www.tapsmart.com/games/golf-battle-fast-paced-online-mini-golf/}. \dostep.

\bibitem{gra-ab} „Angry Birds (video game),” Wikipedia. [Online]. Dostępne:
  \url{https://en.wikipedia.org/wiki/Angry_Birds_(video_game)}. \dostep.

\bibitem{schell} J. Schell, \emph{The Art of Game Design: A~Book of Lenses}, wyd.~3,
  Boca Raton: CRC Press, 2019.

\bibitem{rogers} S. Rogers, \emph{Level Up! The Guide to Great Video Game Design},
  wyd.~2, Chichester: John Wiley \& Sons, 2014.

\bibitem{unity-manual} Unity Technologies, „Unity User Manual 2022.3 (LTS)”. [Online].
  Dostępne: \url{https://docs.unity3d.com/2022.3/Documentation/Manual/index.html}. \dostep.

\bibitem{unity-physics} Unity Technologies, „Introduction to rigid body physics,” Unity
  Manual 2022.3. [Online]. Dostępne:
  \url{https://docs.unity3d.com/2022.3/Documentation/Manual/RigidbodiesOverview.html}.
  \dostep.

\bibitem{unity-raycast} Unity Technologies, „Physics.Raycast,” Unity Scripting API
  2022.3. [Online]. Dostępne:
  \url{https://docs.unity3d.com/2022.3/Documentation/ScriptReference/Physics.Raycast.html}.
  \dostep.

\bibitem{urp} Unity Technologies, „Universal Render Pipeline overview,” URP 14.0.
  [Online]. Dostępne:
  \url{https://docs.unity3d.com/Packages/com.unity.render-pipelines.universal@14.0/manual/index.html}.
  \dostep.

\bibitem{inputsystem} Unity Technologies, „Input System,” dokumentacja pakietu 1.11.
  [Online]. Dostępne:
  \url{https://docs.unity3d.com/Packages/com.unity.inputsystem@1.11/manual/index.html}.
  \dostep.

\bibitem{unity-asmdef} Unity Technologies, „Assembly definitions,” Unity Manual 2022.3.
  [Online]. Dostępne:
  \url{https://docs.unity3d.com/2022.3/Documentation/Manual/ScriptCompilationAssemblyDefinitionFiles.html}.
  \dostep.

\bibitem{seemann} S. van Deursen, M. Seemann, \emph{Dependency Injection Principles,
  Practices, and Patterns}, Shelter Island: Manning, 2019.

\bibitem{nystrom} R. Nystrom, \emph{Game Programming Patterns}, Genever Benning, 2014.
  [Online]. Dostępne: \url{https://gameprogrammingpatterns.com/}. \dostep.

\bibitem{zenject} Modest Tree, „Zenject -- Dependency Injection Framework for Unity3D,”
  repozytorium GitHub. [Online]. Dostępne: \url{https://github.com/modesttree/Zenject}.
  \dostep.

\bibitem{gof} E. Gamma, R. Helm, R. Johnson, J. Vlissides, \emph{Design Patterns: Elements
  of Reusable Object-Oriented Software}, Reading: Addison-Wesley, 1994.

\bibitem{reactivex} ReactiveX, „ReactiveX -- Intro”. [Online]. Dostępne:
  \url{https://reactivex.io/intro.html}. \dostep.

\bibitem{unitask} Cysharp, „UniTask -- Provides an efficient allocation free async/await
  integration for Unity,” repozytorium GitHub. [Online]. Dostępne:
  \url{https://github.com/Cysharp/UniTask}. \dostep.

\bibitem{gossman} J. Gossman, „Introduction to Model/View/ViewModel pattern for building
  WPF apps,” 08.10.2005. [Online]. Dostępne:
  \url{https://learn.microsoft.com/en-us/archive/blogs/johngossman/introduction-to-modelviewviewmodel-pattern-for-building-wpf-apps}.
  \dostep.

\bibitem{ms-mvvm} Microsoft, „Model-View-ViewModel (MVVM),” \emph{Enterprise Application
  Patterns Using .NET MAUI}, 2024. [Online]. Dostępne:
  \url{https://learn.microsoft.com/en-us/dotnet/architecture/maui/mvvm}. \dostep.

\bibitem{nugetforunity} GlitchEnzo, „NuGetForUnity -- A~NuGet Package Manager for Unity,”
  repozytorium GitHub. [Online]. Dostępne:
  \url{https://github.com/GlitchEnzo/NuGetForUnity}. \dostep.

\bibitem{unity-il2cpp} Unity Technologies, „IL2CPP Overview,” Unity Manual 2022.3.
  [Online]. Dostępne: \url{https://docs.unity3d.com/2022.3/Documentation/Manual/IL2CPP.html}.
  \dostep.

\bibitem{android-api} Android Developers, „Android 9 features and APIs”. [Online].
  Dostępne: \url{https://developer.android.com/about/versions/pie/android-9.0}. \dostep.

\bibitem{martin} R. C. Martin, \emph{Czysta architektura. Struktura i~design
  oprogramowania. Przewodnik dla profesjonalistów}, Gliwice: Helion, 2018.

\bibitem{newtonsoft} Newtonsoft, „Json.NET Documentation -- Introduction”. [Online].
  Dostępne: \url{https://www.newtonsoft.com/json/help/html/Introduction.htm}. \dostep.

\bibitem{nystrom-pool} R. Nystrom, „Object Pool,” \emph{Game Programming Patterns}.
  [Online]. Dostępne: \url{https://gameprogrammingpatterns.com/object-pool.html}. \dostep.

\bibitem{cinemachine} Unity Technologies, „About Cinemachine,” dokumentacja pakietu
  Cinemachine 2.10. [Online]. Dostępne:
  \url{https://docs.unity3d.com/Packages/com.unity.cinemachine@2.10/manual/index.html}.
  \dostep.

\bibitem{unity-playerprefs} Unity Technologies, „PlayerPrefs,” Unity Scripting API
  2022.3. [Online]. Dostępne:
  \url{https://docs.unity3d.com/2022.3/Documentation/ScriptReference/PlayerPrefs.html}.
  \dostep.

\bibitem{firebase-unity} Google, „Add Firebase to your Unity project,” dokumentacja
  Firebase. [Online]. Dostępne: \url{https://firebase.google.com/docs/unity/setup}. \dostep.

\bibitem{appodeal} Appodeal, „Get Started,” dokumentacja Appodeal Unity SDK. [Online].
  Dostępne: \url{https://docs.appodeal.com/unity/get-started}. \dostep.

\bibitem{mcconnell} S. McConnell, \emph{Kod doskonały. Jak tworzyć oprogramowanie
  pozbawione błędów}, wyd.~II, Gliwice: Helion, 2010.

\end{thebibliography}
\endgroup
